% Figure : trajectoire evolutive de l'IA generative vers les ecosystemes agentiques
\begin{figure}[H]
\centering
\begin{tikzpicture}[node distance=8mm and 6mm]
  \node[agent] (genai)  {GenAI};
  \node[agent, right=of genai]  (aiagent) {AI Agents};
  \node[agent, right=of aiagent] (agentic) {Agentic AI};
  \node[agent, right=of agentic, text width=2.6cm] (commu) {Communautés agentiques};
  \node[agent, right=of commu, text width=2.8cm] (eco)   {Écosystèmes hybrides gouvernés};

  \draw[flow] (genai)   -- (aiagent);
  \draw[flow] (aiagent) -- (agentic);
  \draw[flow] (agentic) -- (commu);
  \draw[flow] (commu)   -- (eco);

  \node[font=\footnotesize\itshape, below=4mm of aiagent, text width=3cm, align=center] {modules + outils};
  \node[font=\footnotesize\itshape, below=4mm of agentic, text width=3cm, align=center] {planification + mémoire};
  \node[font=\footnotesize\itshape, below=4mm of commu, text width=3cm, align=center] {humains + agents};
  \node[font=\footnotesize\itshape, below=4mm of eco, text width=3cm, align=center] {gouvernance};
\end{tikzpicture}
\caption[Trajectoire évolutive Agentic AI]{Trajectoire évolutive identifiée par la littérature : de l'IA générative aux écosystèmes hybrides gouvernés \cite{schneider2025,milosevic_rabhi}.}
\label{fig:trajectoire}
\end{figure}
